\section{Conditionally Linear Functions, Distributions, and Samplers}
\label{sec:linear}

\subsection{Conditionally linear functions and distributions}
\label{sec:cl}

We first introduce \emph{conditionally linear functions}, which are used to
specify the question distribution for games considered in the paper in a way
that the question distribution can be ``introspected'', as described in
Section~\ref{sec:introspection}.
Intuitively, a conditionally linear function takes as input an element
$x\in V = \F^n$ for some $n\geq 0$, and applies linear maps $L_j$ sequentially
on $x^{V_j}$ where $V_1,V_2,\ldots $ are a sequence of complementary register
subspaces such that both the linear maps $L_j$ and the subspace $V_j$ for
$j \ge 2$ may depend on the values taken by previous linear maps $L_1(x^{V_1})$,
$L_2(x^{V_2})$, etc.

In the remainder of the section we use $V$ to denote the linear space $\F^n$ for
some integer $n \geq 0$.
For ease of notation we extensively use the subscript range notation.
For example, if $V_1, V_2, \ldots, V_\ell$ are fixed subspaces of $V$ and $k \in
\{1, 2, \ldots, \ell\}$ we write
\begin{equation*}
  V_{<k} = \bigoplus_{j\,:\,1 \le j<k} V_j\;,\quad
  V_{> k} = \bigoplus_{j\,:\, \ell\ge j > k} V_j\;,
\end{equation*}
and it is understood that $V_{\le k}$ and $V_{\ge k}$ are identical to
$V_{<k+1}$ and $V_{>k-1}$, respectively.
Moreover, if $V'$ is a register subspace of $V$, $F: V' \to V'$ is a function,
and $x \in V$, we write $x^F$ to denote $F(x^{V'})$.
For example, in the following definition $x^{L_1}$ is used as shorthand notation
for $L_1(x^{V_1})$.


\begin{definition}
  \label{def:cl-func}
	Let $V$ be $\F^n$ for some $n\geq 0$.
  For all integers $\ell\geq 0$ the collection of \emph{$\ell$-level
    conditionally linear functions} (implicitly, \emph{on $V$}) is defined
  inductively as follows.
  \begin{enumerate}
  \item There is a single $0$-level conditionally linear function, which is the
    $0$ function on $V$.

  \item Let $\ell \geq 1$ and suppose the collection of $(\ell-1)$-level
    conditionally linear functions has been defined.
    The collection of $\ell$-level conditionally linear functions on $V$ consists
    of all functions $L$ on $V$ that can be expressed in the following form.
    There exist complementary register subspaces $V_1$ and $V_{>1}$ of $V$, a
    linear function $L_1$ on $V_1$, and for all $v \in L_1(V_1)$, an
    $(\ell-1)$-level conditionally linear function $L_{>1,\, v}$ on $V_{>1}$,
    such that for all $x\in V$,
    \begin{equation*}
      L(x) = x^{L_1} + L_{>1,\, x^{L_1}} (x^{V_{>1}})\;.
    \end{equation*}
  \end{enumerate}
\end{definition}

The concept of conditionally linear functions is simple even though the
notations seem complicated.
It is best understood using an operational definition illustrated in
\cref{fig:cl-functions}.
We give an example of a $2$-level CL function in \cref{ex:cl-function}.

\begin{remark}
  \label{rk:level-1-is-linear}
  It follows from the above definition that all $1$-level CL functions are linear.
\end{remark}

\begin{example}\label{ex:cl-function}
Consider $V = \F_{2}^{3}$. Then the following map
\begin{equation*}
  L:(x_{0}, x_{1}, x_{2})
  \mapsto (0, x_{0}x_{2} + x_{1}x_{2} + x_{0},  x_{2})
\end{equation*}
on $V$ is a $2$-level CL function by choosing $V_{1} = \mathrm{span}\{(0,0,1)\}$, $L_{1}$
the identity function on $V_{1}$, $V_{2} = V_{1}^{\perp}$, and
\begin{equation*}
  L_{2,\, x_{2}}: (x_{0}, x_{1}, 0)
  \mapsto (0, x_{0}x_{2} + x_{1}x_{2} + x_{0}, 0),
\end{equation*}
which is linear for all $x_{2} \in \F_{2}$.
Note that intuitively the function $L$ selects either $x_{0}$ or $x_{1}$ in the
second coordinate conditioning on the value of $x_{2}$.
\end{example}

\begin{figure}[!htbp]
  \centering
  \begin{tikzpicture}



    \node[draw, minimum width=2cm, anchor=east] at (0,0) (xv1) {$x^{V_{1}}$};
    \node[draw, minimum width=6cm, anchor=west] at (-.4pt,0) {$x^{V_{>1}}$};
    \node[draw, minimum width=2cm, anchor=east] at (0,-2) (xl1) {$x^{L_{1}}$};
    \draw[-stealth] (xv1.south) -- ++(0,-1) node[above right] {$L_{1}$} -- (xl1.north);

    \node[draw, minimum width=2cm, anchor=west] at (0,-3) (xv2) {$x^{V_{2}}$};
    \node[draw, minimum width=4cm, anchor=west] at (2cm,-3) {$x^{V_{>2}}$};

    \node[draw, minimum width=2cm, anchor=west] at (0,-5) (xl2) {$x^{L_{2}}$};

    \draw[-stealth] (xv2.south) -- ++(0,-1) node[above right]
      {$L_{2,\,x^{{L_{1}}}}$} -- (xl2.north);
    \draw[dashed, -stealth, to path={|- (\tikztotarget)}] (xl1) edge (xv2.west);

    \draw[dashed, -stealth, to path={|- (\tikztotarget)}] (xl2) edge (3cm,-6);
    \draw[dashed, -stealth, to path={|- (\tikztotarget)}] (xl1) edge (3cm,-6);
    \node[right] at (3cm,-6) {$\;\cdots\cdots$};

  \end{tikzpicture}
  \caption{An illustration of an $\ell$-level CL function $L$.
    First a register subspace $V_{1}$ and a linear map $L_{1}$ are chosen and
    applied to obtain $x^{L_{1}} = L_{1}(x^{V_{1}}) \in V_{1}$.
    Then depending the value of $x^{L_{1}}$, a register subspace $V_{2}$ and a
    linear map $L_{2,\,x^{L_{1}}}$ are chosen and applied to obtain
    $x^{L_{2}} = L_{2,\,x^{{L_{1}}}}(x^{V_{2}}) \in V_{2}$ and so on.
    Finally, $L(x)$ is defiend to be
    $\sum_{k=1}^{\ell} x^{L_{k}}$.}\label{fig:cl-functions}
\end{figure}

\begin{remark}
  \label{rk:higher-level}
  Note that for any integer $\ell\geq 1$ the collection of $\ell$-level CL
  functions trivially contains the collection of $(\ell-1)$-level CL functions:
  for this it suffices to note that the $0$ function, which is a $0$-level CL
  function, is also a $1$-level CL function by setting $V_1=V$, $V_{>1}= \{0\}$,
  $L_1(x)=0$ for all $x\in V$, and $L_{>1,\, x^{L_1}}$ is the $0$ map for all
  $x\in V$.
\end{remark}

\begin{definition}
  \label{def:cl-dist}
  Let $L, R: V \to V$ be conditionally linear functions.
  The \emph{conditionally linear distribution $\mu_{L,R}$ corresponding to $(L,
    R)$} is defined as the distribution over pairs $(L(x), R(x)) \in V \times V$
  for $x$ drawn uniformly at random from $V$.
\end{definition}

Throughout the paper we abbreviate ``conditionally linear functions'' and
``conditionally linear distributions'' as \emph{CL functions} and \emph{CL
  distributions}, respectively.

The following lemma elucidates structural properties of $\ell$-level CL
functions.
Recall that using our shorthand notation, $x^{L_{<k}}$ and $x^{L_k}$ in the
lemma denote $L_{<k}(x)$ and $L_{k,u}(x^{V_{k,\, u}})$ where $u = L_{<k}(x)$.


\begin{lemma}
  \label{lem:cl-kth}
	Let $\ell \geq 1$ and $V = \F^n$ for some integer $n \geq 0$.
  A function $L: V \to V$ is an $\ell$-level CL function if and only if the
  following collection of functions and subspaces exists:
  \begin{enumerate}[label=(\roman*)]
  \item \label{enu:cl-k-marginal} For each $k \in \{1, 2, \ldots, \ell\}$, a
    function $L_{\le k} : V \to V$ called the \emph{$k$-th marginal of $L$};
  \item \label{enu:cl-k-space} For each $k \in \{1, 2, \ldots, \ell\}$ and $u
    \in L_{<k} (V)$, a register subspace $V_{k,\, u}$ of $V$ called the
    \emph{$k$-th factor space with prefix $u$};
  \item \label{enu:cl-k-map} For each $k \in \{1, 2, \ldots, \ell\}$ and $u \in
    L_{<k} (V)$, a linear map $L_{k,\, u} : V_{k,\, u} \to V_{k,\, u}$ called
    the \emph{$k$-th linear map of $L$ with prefix $u$};
  \end{enumerate}
  such that the following conditions hold for all $k \in \{1, 2, \ldots,
  \ell\}$.
  \begin{enumerate}
  \item \label{enu:cl-k-cl} $L_{\leq k}$ is a $k$-level CL function on $V$;
	\item \label{enu:cl-space-sum} $V = \bigoplus_{i = 1}^\ell V_{i,\,
      x^{L_{<i}}}$ for all $x \in V$;
\item \label{enu:cl-map-sum} $L_{\leq k}(x) = \sum_{i = 1}^{k} x^{L_i}$ for
    all $x \in V$, where $L_i$ is shorthand notation for $L_{i,\,
      x^{L_{<i}}}$;
  \item \label{enu:cl-last-ell} $L = L_{\leq \ell}$.
  \end{enumerate}
  As in \cref{enu:cl-map-sum}, we sometimes use $V_k$ and $L_k$ to denote
  $V_{k,\, u}$ and $L_{k,\, u}$ respectively, leaving the prefix $u$ implicit.
\end{lemma}

\begin{proof}
	We first prove the ``if'' direction: if there exist spaces and functions
  satisfying the conditions in the lemma, the fact that $L$ is an $\ell$-level
  CL function follows from \cref{enu:cl-k-cl,enu:cl-last-ell} of the lemma
  statement.

  We now prove the ``only if'' direction.
  Given a CL function $L$ on $V$, we construct the $k$-th family of subspaces
  and functions for all $k \in \{1, \ldots, \ell\}$ by induction on the level
  $\ell$.
  First consider the base case $\ell=1$.
  Since $L_{<1} = 0$, we omit the mentioning of the prefix $u \in L_{<1}(V)$.
  Define $L_{\leq 1} = L$, the factor space $V_1 = V$, and the linear map $L_1 =
  L$.
  It is straightforward to verify that the conditions in the lemma hold for
  these choices of linear maps and spaces.
  
  Now, assume that the lemma holds for CL functions of level at most $\ell - 1$,
  and we prove the lemma for $\ell$-level CL functions.
  By definition, an $\ell$-level CL function $L$ can be written as
  \begin{equation*}
  	L(x) = x^{L_1} + L_{>1,\, x^{L_1}}(x^{V_{>1}})
  \end{equation*}
  for some linear map $L_1 : V_1 \to V_1$ and a family of $(\ell-1)$-level CL
  functions
  \begin{equation*}
    \bigl \{ L_{>1,\, v} : V_{>1} \to V_{>1} \bigr \}_{\, v \in L_1(V_1)}
  \end{equation*}
  where $V_1$ and $V_{>1}$ are complementary register subspaces of $V$.
  Next, using the inductive hypothesis on the $(\ell-1)$-level CL function
  $L_{>1,\, v}$ we get that for all $v \in L_1(V_1)$ and all $k \in \{1, 2,
  \ldots, \ell-1\}$ there exist $k$-th marginal functions $L'_{v,\, \le k}:
  V_{>1} \to V_{>1}$, $k$-th factor spaces $V'_{v,\, k,\, u}$, and $k$-th linear
  maps $L'_{v,\, k,\, u}$ of $L_{>1,\, v}$ with prefix $u \in L'_{v,\,
    <k}(V_{>1})$ such that the conditions of the lemma for $L_{>1,\, v}$ hold.
  
  Define the marginal functions $L_{\le k} : V \to V$, factor spaces
  $V_{k,\, u}$ and linear maps $L_{k,\, u}$ for $L$ as follows.
  \begin{enumerate}[label=(\roman*)]
  \item Define $L_{\le 1} = L_1$ and the first factor space to be $V_1$;
  \item For all $k \in \{2, 3, \ldots, \ell\}$, define
    \begin{equation}
      \label{eq:cl-k-def}
      L_{\le k}: x \mapsto x^{L_1} + L'_{x^{L_1},\, <k} (x^{V_{>1}})
      \quad\text{for } x\in V;
    \end{equation}
	\item For all $k \in \{2, 3, \ldots, \ell\}$ and $u \in L_{< k}(V)$, define
    $V_{k,\, u} = V'_{v,\, k-1,\, w}$ and $L_{k,\, u} = L'_{v,\, k-1,\, w}$
    where $v = u^{V_1} \in L_{1}(V_{1})$ and $w = u^{V_{>1}} \in L'_{v,< k-1}(V_{> 1})$.
  \end{enumerate}
  We verify that the conditions of the lemma are satisfied.
  Since $L'_{v,\, < k}$ is by assumption a $(k-1)$-level CL function on $V_{>
    1}$, we get that $L_{\leq k}$ is a $k$-level CL function on $V$ from
  \cref{eq:cl-k-def}, establishing \cref{enu:cl-k-cl}.
  By the induction hypothesis, we have for all $v\in L_1(V_1)$ and $y \in
  V_{>1}$,
  \begin{equation}
    \label{eq:cl-k-1}
    V_{>1} = \bigoplus_{i=1}^{\ell-1}
    V'_{\smash{v,\, i, \, y^{L'_{v,\, <i}}}},
  \end{equation}
  which implies that for all $x \in V$ and $v = x^{L_1}$,
  \begin{equation*}
    V = V_1 \oplus \Big ( \bigoplus_{i=1}^{\ell-1}
    V'_{\smash{v,\, i, \, x^{L'_{v,\, <i}}}} \Big)
    = V_1 \oplus \Big (\bigoplus_{i=2}^{\ell} V_{i,\, x^{L_{<i}}} \Big)
    = \bigoplus_{i=1}^\ell V_{i,\, x^{L_{<i}}}\;.
  \end{equation*}
  The first equality follows from \cref{eq:cl-k-1} while the second and third
  equalities follow from the definition of $V_{k,\, u}$.
  This establishes \cref{enu:cl-space-sum}.

  Next, we have that for all $x \in V$, $v = x^{L_1}$, and $k \in
  \{1,2,\ldots,\ell\}$,
  \begin{align}
    \label{eq:cl-last-ell}
    L_{\leq k}(x)
    & = v + L'_{v,\, <k} (x^{V_{>1}}) \\
    & = v + \sum_{i=1}^{k-1} x^{L'_{v,\, i}}
    = \sum_{i = 1}^{k} x^{L_{i}}, \label{eq:cl-last-ell2}
  \end{align}
  where $L'_{v,\, i}$ is the $i$-th linear map of $L_{>1,\, v}$ with prefix
  $L'_{v,\, <i}(x)$ and $L_i$ is the $i$-th linear map of $L$ with prefix
  $L_{<i}(x)$.
  The second equality follows from the inductive hypothesis applied to $L'_{v,\,
    <k}$ and the third equality follows from the definition of $L_{i}$.
  Line~\eqref{eq:cl-last-ell2} implies \cref{enu:cl-map-sum} of the lemma.
  
  Finally, \Cref{enu:cl-last-ell} follows from~\eqref{eq:cl-last-ell} where we
  set $k = \ell$ and observe that $L'_{x^{L_1},\, \le \ell-1}$ is equal to
  $L_{>1,\, x^{L_1}}$ by the inductive hypothesis.
  This shows that $L_{\le k}$, $V_{k,\, u}$, and $L_{k,\, u}$ satisfy the
  conditions of the lemma and completes the induction.
\end{proof}

We note that the marginal functions, factor spaces, and linear maps of a given
CL function $L$ may not be unique; for example, consider the identity function
on a linear space $V = \F^n$.
This is clearly a $1$-level CL function, but it can also be viewed as a
$k$-level CL function for $k \in \{2, \ldots, n\}$ with an arbitrary partition
of $V$ into factor spaces.

\begin{lemma}
  \label{lem:cl-concat}
  Let $\ell, k \ge 0$ be integers and $U = \F^n$, $V = \F^m$ be linear spaces.
  Suppose $L$ is a $k$-level CL function on $U$ and $R_u$ is an $\ell$-level CL
  function on $V$ for each $u \in L(U)$.
  Then the \emph{concatenation} $T$ of $L$ and $\{R_u\}_u$ defined as
  \begin{equation*}
    T (x) = L(x^U) + R_{L(x^U)} (x^V)
  \end{equation*}
  is a $(k+\ell)$-level conditionally linear function on $U\oplus V$.
\end{lemma}

\begin{proof}
  We prove the claim by induction on $k$.
  The case $k = 0$ follows from the Definition~\ref{def:cl-func}.
  Assume that the lemma holds for $L$ being at most $(k-1)$-level.
  By Definition~\ref{def:cl-func}, there are complementary register subspaces
  $U_1$ and $U_{>1}$ of $U$, a linear function $L_1$ on $U_1$, and a family of
  $(k-1)$-level CL functions $L_{>1,\, v}$ for $v \in L_1(U_1)$ such that
  \begin{equation*}
    L(x^U) = x^{L_1} + L_{>1,\, x^{L_1}}(x^{U_{>1}}).
  \end{equation*}
  For all $x^{L_1}$, define the function $T_{>1,\, x^{L_1}}$ on $U_{>1} \oplus V$ as
  \begin{equation*}
    T_{>1,\, x^{L_1}}(x^{U_{>1} \oplus V}) = L_{>1,\, x^{L_1}}(x^{U_{>1}}) +
    R_{x^{L_1} + x^{L_{>1}}}(x^V),
  \end{equation*}
  the concatenation of $L_{>1,\, x^{L_1}}$ and $\{ R_{L(x^U)} \}_{x^{L_{>1}}}$
  where $L_{>1}$ is the shorthand notation of $L_{>1,\, x^{L_1}}$.
  By the induction hypothesis, $T_{>1,\, x^{L_1}}$ is $(k+\ell-1)$-level
  conditionally linear.
  The lemma follows from Definition~\ref{def:cl-func}.
\end{proof}
\begin{lemma}[Direct sums of CL functions]
  \label{lem:cl-func-prod}
  Let $V^{(1)}, V^{(2)}, \ldots, V^{(m)}$ be register subspaces of $V$ such that
  $V = \bigoplus_{j=1}^m V^{(j)}$.
  Suppose that, for each $j\in\{1, 2, \ldots, m\}$, $L^{(j)}$ is an
  $\ell_j$-level conditionally linear function on $V^{(j)}$.
  Then the direct sum $L = \bigoplus_{j=1}^m L^{(j)}$ is an $\ell$-level CL
  function over $V$ for $\ell = \max_j \{\ell_j\}$, where $L$ is defined by
  \begin{equation*}
    L(x) = \sum_{j=1}^m  L^{(j)}(x^{(j)})
  \end{equation*}
  for all $x = \sum_{j=1}^m x^{(j)} \in \bigoplus_{j=1}^m V^{(j)}$.
\end{lemma}

\begin{proof}
  It is easy to see that an $\ell$-level CL function is also $k$-level
  conditionally linear for all $k \ge \ell$.
  Hence, it suffices to prove the claim where $\ell_j = \ell$ for $j=1, 2,
  \ldots, m$.

  We prove the theorem by an induction on $\ell$.
  For $\ell = 1$, the functions $L^{(j)}$ are linear and the claim follows by
  the fact that the direct sum of linear maps is linear.

  Assume now the theorem holds for conditionally linear functions of level at
  most $\ell-1$ and $L^{(j)}$ are $\ell$-level conditionally linear functions
  for $j=1, 2, ,\ldots, m$.
  By definition, $L^{(j)}$ is the concatenation of conditionally linear
  functions $L^{(j)}_{1}$ on $V^{(j)}_1$ and $\{ L^{(j)}_{>1,\, v_j} \}_{v_j}$
  on $V^{(j)}_{>1}$ of levels $1$, and $\ell-1$ respectively.
  Furthermore,
  \begin{equation*}
    L(x) = \sum_{j=1}^m L^{(j)} (x^{(j)}) = \sum_{j=1}^m \Bigl( v_j +
    L^{(j)}_{>1,\, v_j}
    \Bigl((x^{(j)})^{V^{(j)}_{>1}} \Bigr) \Bigr),
  \end{equation*}
  where $v_j = L^{(j)}_{1}\Bigl( (x^{(j)})^{V^{(j)}_1}
    \Bigr)$.
  By the induction hypothesis,
  \begin{equation*}
    L_{1} (x^{V_1}) = \sum_{j=1}^m L^{(j)}_{1}\Bigl( (x^{(j)})^{V^{(j)}_1}
    \Bigr),\quad L_{>1,\, v} (x^{V_{>1}}) = \sum_{j=1}^m L^{(j)}_{>1,\, v_j}
    \Bigl((x^{(j)})^{V^{(j)}_{>1}} \Bigr)
  \end{equation*}
  are $1$-level and $(\ell-1)$-level conditionally linear respectively for $v =
  \sum_j v_j$, $V_1 = \bigoplus_{j=1}^m V^{(j)}_1$, and $V_{>1} =
  \bigoplus_{j=1}^m V^{(j)}_{>1}$.
  This proves that $L$ is $\ell$-level conditionally linear.
\end{proof}

\begin{lemma}
  \label{lem:cl-dist-prod}
  For each $i \in \{1,\ldots,m\}$ let $L^{(i)}, R^{(i)} : V^{(i)} \to V^{(i)}$
  be $\ell_i$-level conditionally linear functions and let $L, R: V \to V$ be the
  direct sums $L = \bigoplus_i L_i$ and $R = \bigoplus_i R_i$, respectively, as
  defined in Lemma~\ref{lem:cl-func-prod}.
  Then the conditionally linear distribution $\mu_{L, R}$ is the product
  distribution $\prod_{i=1}^m \mu_{L^{(i)},\, R^{(i)}}$ over $V \times V$.
\end{lemma}

\begin{proof}
	The distribution $\mu_{L,R}$ is the distribution over pairs $(L(x),R(x))$
  where $x$ is sampled uniformly from $V \times V$.
  By Lemma~\ref{lem:cl-func-prod}, this is equivalent to the distribution over
  pairs $\bigl( (L_i(x^{V_i}))_{i=1}^m, (R_i(x^{V_i}))_{i=1}^m \bigr)$ where $x$
  is chosen uniformly at random from $V$.
  This distribution is exactly the product of the distributions $
  \mu_{L^{(i)},\, R^{(i)}}$ for $i = 1, 2, \ldots, m$.
\end{proof}

CL functions used in the paper are frequently defined over a ``large'' field
$\F_{2^t}$ (e.g., the CL functions used in the low degree tests of
Section~\ref{sec:ldt}).
However, the introspection protocol in Section~\ref{sec:introspection} handles
CL functions defined over $\F_2$.
The following definition and lemma show that CL functions over prime power
fields can be viewed as CL functions over the prime field via a ``downsizing''
operation.

\begin{definition}[Downsizing CL functions]
  \label{def:cl-downsize}
  Let $V = \F_{q}^n$ be a linear space for a prime power $q=p^t$ for odd $t$.
  Let $L: V \to V$ be a function.
  Let $\downsize(\cdot)$ denote the downsize map from
  Section~\ref{sec:finite-fields} corresponding to the basis
  $\{e_1,\ldots,e_t\}$ of $\F_q$ over $\F_p$ specified by
  Lemma~\ref{lem:efficient_basis}.
  In particular, $\downsize$ is linear over $\F_p$, and by
  Lemma~\ref{lem:downsize_subspace}, the set $\downsize(V)$ is the linear space
  $\F_p^{nt}$.
  Define the \emph{downsized function} $L^\downsize : \downsize(V) \to
  \downsize(V)$ by $L^\downsize = \downsize \circ L \circ \downsize^{-1}$.
\end{definition}

\begin{lemma}
  \label{lem:cl-downsize}
	Let $V = \F_q^n$ for a prime power $q = p^t$.
  Let $L: V \to V$ be an $\ell$-level CL function over $V$ for some integer
  $\ell \ge 0$.
  Let $L_{\leq j}$, $V_{j,\, u}$, and $L_{j,\, u}$ denote the $j$-th marginal
  functions, factor spaces, and linear maps corresponding to $L$ as guaranteed
  by Lemma~\ref{lem:cl-kth}.
  Then $L^\downsize: \downsize(V) \to \downsize(V)$ is an $\ell$-level CL
  function on $V^\downsize = \downsize(V) = \F_p^{nt}$ with marginal functions
  $L^\downsize_{\leq j}$, factor spaces $V^\downsize_{j,\, v}$, and linear maps
  $L^\downsize_{j,\, v}$ that satisfy the following for all
  $j\in\{1,\ldots,\ell\}$.
  \begin{enumerate}
  \item The $j$-th marginal function $L^\downsize_{\leq j}$ of $L^\downsize$ is
    equal to $\downsize \circ L_{\leq j} \circ \downsize^{-1}$.
	\item For all $u \in L_{< j} (V)$, the $j$-th factor space $V_{j,\,
      \downsize(u)}^\downsize$ and the $j$-th linear map $L_{j,\,
      \downsize(u)}^\downsize$ of $L^\downsize$ are equal to $\downsize(V_{j,\,
      u})$ and $\downsize \circ L_{j,\, \downsize(u)} \circ \downsize^{-1}$ respectively.
  \end{enumerate}
\end{lemma}

\begin{proof}
  We prove the lemma by induction on $\ell$.
  Let $L:V \to V$ be an $\ell$-level CL function.
  For the base case $\ell = 1$, observe that since $\downsize$ is a linear
  bijection between $\F_q$ and $\F_p^t$ as linear spaces over $\F_p$, the
  function $L^\downsize$ is linear, and thus a $1$-level CL function over
  $\downsize(V) = \F_p^{nt}$.
  Furthermore, the first marginal function $L^\downsize_{\leq 1} = L^\downsize =
  \downsize \circ L_{\leq 1} \circ \downsize^{-1}$; the factor space
  $V^\downsize_1 = \downsize(V_1) = \downsize(V)$, and $L^\downsize_1 = L =
  \downsize \circ L_1 \circ \downsize^{-1}$.

  Assume that the statement of the lemma holds for some $\ell - 1 \geq 1$.
  Let $L_{\leq j}$, $V_{j,\, u}$, and $L_{j,\, u}$ denote the marginal
  functions, factor spaces, and linear maps corresponding to $L$ as guaranteed
  by Lemma~\ref{lem:cl-kth}.
  Recursively define the following functions and spaces, for $j \in
  \{1,\ldots,\ell\}$.
  \begin{enumerate}
	\item $L^\downsize_{\leq j} = \downsize \circ L_{\leq j} \circ
    \downsize^{-1}$.
	\item For all $u \in L_{<j} (V)$, set $V_{j,\, \downsize(u)}^\downsize =
    \downsize(V_{j,\, u})$ and set $L_{j,\, \downsize(u)}^\downsize = \downsize
    \circ L_{j,\,\downsize( u)} \circ \downsize^{-1}$.
  \end{enumerate}
  We argue that $\{L^\downsize_{\leq j}\}$, $\{V^\downsize_{j,\, v}\}$, and
  $\{L_{j,\, v}^\downsize \}$ satisfy the conditions of Lemma~\ref{lem:cl-kth}
  for the function $L^\downsize$, which implies that $L^\downsize$ is an
  $\ell$-level CL function over $\downsize(V)$.

  We first establish \Cref{enu:cl-last-ell} of Lemma~\ref{lem:cl-kth}.
  Since $L_{\leq \ell} = L$, this implies
  \begin{equation*}
	  L^\downsize = \downsize \circ L \circ \downsize^{-1}
    = \downsize \circ L_{\leq \ell} \circ \downsize^{-1}
    = L^{\downsize}_{\leq \ell},
  \end{equation*}
  as desired.
  Next, for all $j \in \{1, 2, \ldots, \ell\}$, for all $y \in \downsize(V)$
  with $y = \downsize(x)$ for some $x \in V$, letting $u_j = L_{<j}(x)$, we have
  \begin{equation*}
	  \downsize(V) = \downsize \biggl ( \bigoplus_{j=1}^\ell V_{j,\, \smash{x^{L_{<j}}}} \biggr)
    = \bigoplus_{j=1}^\ell \downsize \bigl( V_{j,\, \smash{x^{L_{<j}}}} \bigr)
    = \bigoplus_{j=1}^\ell V^{\,\downsize}_{j,\, \smash{\downsize(x^{L_{<j}})}}\,.
  \end{equation*}
  The first equality follows from \cref{enu:cl-space-sum} of \cref{lem:cl-kth},
  the second equality follows from \cref{lem:downsize_subspace}, and the third
  equality follows by definition.
  Since $\downsize(x^{L_{<j}}) = L^\downsize_{<j} \circ \downsize (x) =
  L^\downsize_{<j} (y)$, this establishes \cref{enu:cl-space-sum} of
  \cref{lem:cl-kth}.

  Next, we have for all $j \in \{1, 2, \ldots, \ell\}$ and all $y \in
  \downsize(V)$ with $y = \downsize(x)$ for some $x \in V$,
  \begin{equation*}
    \begin{split}
      L_{\leq j}^\downsize(y)
      & = \bigl( \downsize \circ L_{\leq j} \circ \downsize^{-1} \bigr) (y)
      = \bigl( \downsize \circ L_{\le j} \bigr) (x)
      = \sum_{i=1}^{j} \downsize \bigl( x^{L_{i,\, x^{L_{<i}}}} \bigr)\\
      & = \sum_{i=1}^{j} L^\downsize_{i,\, v_i}
      \bigl( \downsize(x^{V_{i,\, x^{L_{<i}}}}) \bigr)
      = \sum_{i=1}^{j} L^\downsize_{i,\, v_i} (y^{V^\downsize_{i,\, v_i}})
    \end{split}
  \end{equation*}
  where $v_i = L^\downsize_{< i}(y) = \downsize(x^{L_{<i}})$.
  The first equality follows from definition of $L^\downsize_{\leq j}$, the
  second equality follows from $y = \downsize(x)$, the third equality follows
  from \cref{enu:cl-map-sum} of \cref{lem:cl-kth} applied to $L_{\leq j}$, the
  fourth equality follows from the definition of the linear map
  $L^\downsize_{i,\, v}$, and the fifth equality follows from
  Lemma~\ref{lem:downsize_subspace}.
  This establishes \cref{enu:cl-map-sum} of Lemma~\ref{lem:cl-kth} for
  $L^\downsize_{\leq j}$.

  Finally, since $L_{\leq j}$ is a $j$-level CL function over $V$, using the
  inductive hypothesis we have that $L^\downsize_{\leq j}$ is a $j$-level CL
  function over $\downsize(V)$ when $j \in \{1,2,\ldots,\ell-1\}$.
  It remains to establish that $L_{\leq \ell}^\downsize$ is an $\ell$-level CL
  function.
  Since $L$ is an $\ell$-level CL function, there exists register subspaces
  $V_1, V_{> 1}$ such that $V = V_1 \oplus V_{> 1}$, a linear map $L_1 : V_1 \to
  V_1$ and a collection of $(\ell-1)$-level CL functions $\{ L_{>1,\, v} : V_{>
    1} \to V_{> 1} \}_{v \in L_1(V_1)}$ such that $L(x) = x^{L_1} + L_{>1,\,
    x^{L_1}}(x^{V_{>1}})$ for all $x \in V$.
  Observe that $L_1^\downsize$ is a $1$-level CL function on $V_1^\downsize =
  \downsize(V_1)$, and for $v' = \downsize(v) \in L_1^\downsize(V_1^\downsize)$,
  the inductive hypothesis implies the function $L_{>1,\, v'}^\downsize$ is an
  $(\ell-1)$-level CL function on $V_{>1}^\downsize = \downsize(V_{>1})$.
  Furthermore, since $L^\downsize_{\leq \ell} = L^\downsize$, we have that for
  all $y \in \downsize(V)$ with $y = \downsize(x)$ for some $x \in V$,
  \begin{gather*}
  	L^\downsize_{\leq \ell}(y) = L^\downsize(y) = L_1^\downsize(y) +
    L^\downsize_{>1,\, L_1^\downsize(y)}(y^{V_{>1}^\downsize})
  \end{gather*} 
  which implies that $L^\downsize_{\leq \ell}$ is an $\ell$-level CL function
  over $\downsize(V) = \downsize(V_1) \oplus \downsize(V_{>1})$.
  This establishes \Cref{enu:cl-k-cl} of Lemma~\ref{lem:cl-kth}, and completes
  the induction.
\end{proof}

\begin{lemma}\label{lem:downsize-cl-dist}
  Let $V = \F_q^n$ for some integer $n$ and prime power $q=p^t$ for odd $t$.
  Let $L,R : V\to V$ be CL functions.
  Let $L^\downsize, R^\downsize : \downsize(V)\to \downsize(V)$ be the
  associated downsized CL functions, as defined in
  Definition~\ref{def:cl-downsize}.
  Then the distribution $\mu_{L^\downsize,R^\downsize}$ over $\downsize(V)\times
  \downsize(V)$ defined in Definition~\ref{def:cl-dist} is identical to the
  distribution of $(x,y) \in \downsize(V)\times \downsize(V)$ obtained by first
  sampling $(x',y')$ according to $\mu_{L,R}$ and then returning
  $(\downsize(x'),\downsize(y'))$.
\end{lemma}

\begin{proof}
  The fact that $L^\downsize, R^\downsize$ are well-defined CL functions follows
  from Lemma~\ref{lem:cl-downsize}.
  The lemma is immediate from the definition of $\mu_{L^\downsize,R^\downsize}$
  and the fact that $\downsize$ is a bijection.
\end{proof}

\subsection{Conditionally linear samplers}
\label{sec:cls}

Samplers are Turing machines that perform computations corresponding to CL
functions defined in Section~\ref{sec:cl}.
The inputs and outputs of the sampler are binary strings that are interpreted as
representing data of different types (integers, bits, vectors in $\F_q^s$,
etc.).
See Section~\ref{sec:ff-representations} and in particular
Remark~\ref{rmk:tm_fields} for an in-depth discussion of representing structured
objects on a Turing machine.

\begin{definition}
	A function $q: \N \to \N$ is an \emph{admissible field size function} if for
  all $n \in \N$, $q(n)$ is an admissible field size as defined in
  Definition~\ref{def:admissible-size}.
\end{definition}


\begin{definition}[Conditionally linear samplers]
  \label{def:sampler}
	Let $q:\N \to \N$ be an admissible field size function, and let $s: \N \to \N$
  be a function.
  A $6$-input Turing machine $\sampler$ is a \emph{$\ell$-level conditionally
    linear sampler with field size $q(n)$ and dimension $s(n)$} if for all $n
  \in \N$, letting $q = q(n)$ and $s = s(n)$, there exist $\ell$-level CL
  functions $L^{\alice,\, n}, L^{\bob,\, n} : \F_q^s \to \F_q^s$ with marginal functions
  $\{ L^{w,\, n}_{\leq j}\}$ and factor spaces $\{V^{w,\, n}_{j,\, u}\}$ for $w \in \AB$
  satisfying the conditions of Lemma~\ref{lem:cl-kth}, such that for all $w \in
  \AB$, $j \in \{1,\ldots,\ell\}$, $z \in \F_q^s$:
  \begin{itemize}
  \item On input $(n, \gamestyle{dimension})$, the sampler $\sampler$ outputs the
    dimension $s(n)$.
  \item On input $(n, w, \gamestyle{marginal}, j, z)$, the sampler $\sampler$
    outputs the binary representation of $L^{w,\, n}_{\leq j}(z)$,
	 \item On input $(n, w, \gamestyle{linear}, j, u, y)$, the sampler $\sampler$
    outputs the binary representation of $L^{w,\, n}_{j,\, u}(y)$, where $u$ is
    interpreted as an element of $V^{w,\, n}_{<j}$,
	\item On input $(n, w, \gamestyle{factor}, j, u)$, the sampler $\sampler$
    outputs the $j$-th factor space $V^{w,\, n}_{j,\, u}$ of $L^{w,\, n}$ with
    prefix $u \in L^{w,\, n}_{<j}(V)$, represented as a vector in
    $\{0,1\}^s$ indicating which elementary basis vectors of $\F_q^s$ span the factor space.
  \end{itemize}
  We call $\F_{q(n)}^{s(n)}$ the \emph{ambient space of $\sampler$} on index $n$.
  We call the CL functions $L^{w,\, n}$ for $w \in \ab$ the \emph{CL functions
    of $\sampler$ on index $n$}.
  The \emph{time complexity} of $\sampler$, denoted as $\TIME_\sampler(n)$, is
  the number of steps before $\sampler$ halts for index $n$.
\end{definition}

\begin{remark}
  \label{rmk:sampler-inputs}
  Conditionally linear samplers are defined to have $6$-input tapes, but
  depending on the input, not all input tapes are read.
  For example, if the second input tape has the input $\gamestyle{dimension}$,
  then the remaining input tapes are ignored.
  Thus for notational convenience we write samplers with different numbers of
  arguments, depending on the type of argument it gets.
  The number of arguments is always at most $6$, however.
\end{remark}

The following definition shows how samplers naturally correspond to
conditionally linear distributions.
\begin{definition}[Distribution of a sampler]
  \label{def:sampler-sample}
  Let $\sampler$ be a sampler with field size $q(n)$, dimension $s(n)$.
  For each $n \in \N$, let $L^{\alice,\, n}, L^{\bob,\, n}$ denote the CL
  functions of $\sampler$ on index $n$.
  Let $\mu_{\sampler,\, n}$ denote the CL distribution $\mu_{L^{\alice, n},\,
    L^{\bob, n}}$ corresponding to $(L^{\alice,\, n}, L^{\bob,\, n})$, as
  defined in Definition~\ref{def:cl-dist}.
  We call $\mu_{\sampler,\, n}$ the \emph{distribution of sampler $\sampler$ on
    index $n$}.
\end{definition}

The following provides a definition of a ``downsized'' sampler that can be
obtained from any sampler $\sampler$ over an admissible field $\F_q$.
\begin{definition}[Downsized sampler]
  \label{def:downsize_sampler}
	Let $q: \N \to \N$ be an admissible field size function.
  Let $\sampler$ be an $\ell$-level sampler with field size $q(n)$ and dimension
  $s(n)$.
  Define $\downsize(\sampler)$ as the following Turing machine.
  For all $n \in \N$, $w \in \{\alice,\bob\}$, $j \in \{1,\ldots,\ell\}$, and $z
  \in \F_2^{s \log q}$ where $q = q(n)$ and $s = s(n)$:
	\begin{itemize}
  \item On input $(n, \gamestyle{dimension})$, the sampler returns the output of
    $\sampler(n, \gamestyle{dimension})$ multiplied by $\log q$.
		
	\item On input $(n, w, \gamestyle{marginal}, j, z)$, the sampler
    $\downsize(\sampler)$ returns the output of $\sampler(n, w,
    \gamestyle{marginal}, j, z)$.
		
  \item On input $(n, w, \gamestyle{linear}, j, u', y)$, the sampler
    $\downsize(\sampler)$ computes $u$ such that $u' = \downsize(u)$ and returns
    the output of $\sampler(n, w, \gamestyle{linear}, j, u, y)$.
		
	\item On input $(n, w, \gamestyle{factor}, j, u')$, the sampler
    $\downsize(\sampler)$ computes $u$ such that $u' = \downsize(u)$ and the
    indicator vector
    \[
      C = \sampler(n, w, \gamestyle{factor}, j, u) \in \{0,1\}^s\;,
    \]
    and returns the expanded indicator vector $(D_1, D_2, \ldots, D_s) \in (\{0,
    1\}^{\log q})^s$ where $D_i$ is the all ones vector in $\{0, 1\}^{\log q}$
    if $C_i = 1$ and $D_i$ is the all zeroes vector otherwise.
  \end{itemize}
\end{definition}

The next lemma establishes that $\downsize(\sampler)$ is a well-defined CL
sampler, in the sense that it can be derived from a family of CL functions as in
Definition~\ref{def:sampler}.

\begin{lemma}
  \label{lem:downsize_sampler}
  Let $\ell \geq 1$ be such that $\sampler$ is an $\ell$-level CL sampler, and
  let $q(n)$ and $s(n)$ be as in Definition~\ref{def:downsize_sampler}.
  Then the Turing machine $\downsize(\sampler)$ is an $\ell$-level CL sampler
  with field size $2$, dimension $s'(n) = s(n) \log q(n)$, and
  time complexity
  \begin{equation*}
\TIME_{\downsize(\sampler)}(n) = O \bigl( \TIME_\sampler(n) \log q(n) \bigr)\;.
  \end{equation*}
  Furthermore, for every integer $n \in \N$ the CL functions of
  $\downsize(\sampler)$ on index $n$ are $(L^{\alice,\, n})^\downsize$ and
  $(L^{\bob,\, n})^\downsize$, where $L^{\alice,\, n}, L^{\bob,\, n}$ are the CL
  functions of $\sampler$ on index $n$.
\end{lemma}

\begin{proof} 
  To show that $\downsize(\sampler)$ is an $\ell$-level CL sampler we first show
  the ``Furthermore'' part, i.e.\ verify that for any integer $n\geq 1$ the CL
  functions $(L^{\alice,\, n})^\downsize$ and $(L^{\bob,\, n})^\downsize$ are its
  associated CL functions on index $n$, as defined in
  Definition~\ref{def:sampler}.

  Observe that for $z\in V$, the binary representation of $z$ as an element of
  $\{0,1\}^{s\log q}$ passed as input to $\sampler$ is, by definition (see
  Section~\ref{sec:ff-representations}), identical to the binary representation
  of $\downsize(z)$.
  Using the definition $(L^{w,\, n})^\downsize = \downsize \circ L^{w,\, n}
  \circ \downsize^{-1}$ for $w \in \ab$ this justifies that
  $\downsize(\sampler)$ returns the correct output when executed on inputs of
  the form $(n,\gamestyle{dimension})$, $(n,w,\gamestyle{marginal},j,z)$ and
  $(n, w, \gamestyle{linear}, j, u', y)$.

  Next, if $T$ is a register subspace of $\F_q^s$ with indicator vector $C \in
  \{0,1\}^s$, then $\downsize(T)$ is a register subspace of $\F_2^{s\log q}$
  with indicator vector $D$ defined from $C$ as in
  Definition~\ref{def:downsize_sampler}.
  Thus the output of $\downsize(\sampler)$ on input $(n, w, \gamestyle{factor},
  j, u')$ is equal to the indicator vector of $\downsize(V^{w,\, n}_{j,\, u})$,
  which is the $j$-th factor space of $L^{w,\, n}$ with prefix $u' =
  \downsize(u)$.

  The time complexity of $\downsize(\sampler)$ is the same as
  with the sampler $\sampler$, except it takes $O(\log q(n))$ times longer to
  output the factor space indicator vectors.

\end{proof}


 
\section{Nonlocal Games and $\MIP^*$}
